\section{Context Assembly and Prompt Construction}
\label{sec:gr-context-assembly}

Once a relevant subgraph has been retrieved, it must be serialized into a form the
language model can consume, since the generator ultimately conditions on a textual
prompt~\cite{lewis2020rag}. In Graph RAG this means linearizing entities,
relations, and any attached community summaries into the context window, so that
the structure retrieved from the graph is preserved as far as possible in the
prompt~\cite{peng2024graphrag}. The design of this step matters because the answer
is conditioned jointly on the query and the assembled evidence; poorly organized
context degrades the generation even when the right facts were
retrieved~\cite{gao2023ragsurvey}.

Because context windows are finite, assembly also involves selection and ordering:
the GraphRAG global pipeline, for instance, first generates partial answers from
individual community summaries and only then combines them, rather than attempting
to fit an entire corpus into one prompt~\cite{edge2024graphrag}. Dual-level
retrieval schemes similarly assemble both specific entities and broader relational
context so that the prompt reflects both local detail and global
theme~\cite{guo2024lightrag}.
